\begin{minipage}{\linewidth}
\noindent\textbf{1.}  -- \textbf {(B).} When adjusting your mirrors to carry a passenger, the passenger should:  \textbf {Sit on the seat with you.} \par \textit {A passenger's weight changes the bike's attitude and your sitting position, so mirrors must be set with the passenger seated, exactly as you will ride.} \\
\end{minipage}

\begin{minipage}{\linewidth}
\noindent\textbf{2.}  -- \textbf {(C).} A passenger should:  \textbf {Hold onto the operator's waist.} \par \textit {Holding the operator's waist (or hips/belt) lets the passenger synconiusly move with the rider and lean with the bike, keeping the motorcycle stable.} \\
\end{minipage}

\begin{minipage}{\linewidth}
\noindent\textbf{3.}  -- \textbf {(A).} When the motorcycle is stopped at an intersection, a passenger should:  \textbf {Keep both feet on the foot pegs.} \par \textit {Putting feet down can unbalance the bike and risks injury, so only the operator handles the stop.} \\
\end{minipage}

\begin{minipage}{\linewidth}
\noindent\textbf{4.}  -- \textbf {(C).} When carrying a passenger:  \textbf {It will take longer to accelerate and slow down.} \par \textit {Even a small passenger is a big share of a motorcycle's total weight, far more than two or three passengers are in a car. This makes the bike harder to handle and slower to accelerate and stop.} \\
\end{minipage}

\begin{minipage}{\linewidth}
\noindent\textbf{5.}  -- \textbf {(A).} When riding with passengers, you should:  \textbf {Start slowing sooner.} \par \textit {Even a small passenger is a big share of a motorcycle's total weight, far more than two or three passengers are in a car. This makes the bike harder to handle and slower to accelerate and stop.} \\
\end{minipage}

\begin{minipage}{\linewidth}
\noindent\textbf{6.}  -- \textbf {(D).} When carrying loads:  \textbf {Keep the load low and as far forward as possible.} \par \textit {A low, forward load keeps bike's original center of gravity unchanged, preserving stability and handling. A load behind the rear axle can cause wobble or weave.} \\
\end{minipage}

\begin{minipage}{\linewidth}
\noindent\textbf{7.}  -- \textbf {(D).} A group of riders move from a staggered formation into a single file when:  \textbf {Turning corners.} \par \textit {Turns, curves, and entering or leaving a highway require a rider to have full lane space for room to react and maneuver, so the group goes single file. And then regroup to a staggered formation again.} \\
\end{minipage}

\begin{minipage}{\linewidth}
\noindent\textbf{8.}  -- \textbf {(D).} In this picture, riders are in a staggered formation. The following distance between A and C should be at least:  \textbf {Two seconds.} \par \textit {A staggered formation keeps the group compact while each rider still keeps a two-second gap to the bike directly ahead, giving room to react and maneuver.} \\
\end{minipage}

\begin{minipage}{\linewidth}
\noindent\textbf{9.}  -- \textbf {(C).} Riders in a staggered formation will be passing a car. After the lead rider passes, he/she should:  \textbf {Return to the left portion of the lane.} \par \textit {Riders pass one at a time, even in a staggered formation. After passing, the lead rider returns to the left portion of the lane, leaving room for the next rider, who moves up as the new lead and passes the same way, until the whole group is past.} \\
\end{minipage}

\begin{minipage}{\linewidth}
\noindent\textbf{10.}  -- \textbf {(A).} Alcohol enters the bloodstream and quickly reaches the brain. What is affected first?  \textbf {Judgment.} \par \textit {Alcohol impairs judgment and inhibitions first, so a rider may take risks and misjudge their own impairment before physical effects like slurred speech, balance, or vision problems show up.} \\
\end{minipage}

\begin{minipage}{\linewidth}
\noindent\textbf{11.}  -- \textbf {(C).} Alcohol first affects your:  \textbf {Judgment.} \par \textit {Alcohol impairs judgment and inhibitions first, so a rider may take risks and misjudge their own impairment before physical effects like slurred speech, balance, or vision problems show up.} \\
\end{minipage}

\begin{minipage}{\linewidth}
\noindent\textbf{12.}  -- \textbf {(C).} Approximately half of all motorcycle crash fatalities are caused by:  \textbf {Alcohol Impairment.} \par \textit {Alcohol impairs judgment and inhibitions, so a rider may take risks and misjudge their own impairment before physical effects like slurred speech, balance, or vision problems show up.} \\
\end{minipage}

\begin{minipage}{\linewidth}
\noindent\textbf{13.}  -- \textbf {(B).} If you don't control your drinking:  \textbf {Allow at least an hour for each drink before riding.} \par \textit {Only time removes alcohol, at roughly one drink per hour. Riding slower, exercise, or coffee don't lower impairment.} \\
\end{minipage}

\begin{minipage}{\linewidth}
\noindent\textbf{14.}  -- \textbf {(B).} Studies show that almost half of all riders killed in motorcycle crashes:  \textbf {Had been drinking.} \par \textit {Alcohol is a factor in a large share of fatal motorcycle crashes. Don't drink and ride.} \\
\end{minipage}

\begin{minipage}{\linewidth}
\noindent\textbf{15.}  -- \textbf {(A).} The ability to think clearly and ride safely is affected by as little as:  \textbf {1 drink.} \par \textit {Even one drink can slow reactions and impair judgment, so there is no safe amount for a rider. Don't drink and ride.} \\
\end{minipage}

\begin{minipage}{\linewidth}
\noindent\textbf{16.}  -- \textbf {(A).} If you drink 4 typical drinks, what is the least amount of alcohol remaining in your body after 2 hrs?  \textbf {2 drinks.} \par \textit {The body clears about one drink per hour, so after 2 hours at most 2 of the 4 drinks are gone, leaving at least 2 in your system.} \\
\end{minipage}

\begin{minipage}{\linewidth}
\noindent\textbf{17.}  -- \textbf {(C).} For greatest safety, your helmet should:  \textbf {Fit snugly all the way around.} \par \textit {A snug, even fit keeps the helmet from shifting or coming off in a crash. Helmets save lives. Wear full face helmet. Don’t tempt fate. Don't test the waters.} \\
\end{minipage}

\begin{minipage}{\linewidth}
\noindent\textbf{18.}  -- \textbf {(B).} The single most important thing you can do to improve your chances of surviving a crash is to:  \textbf {Wear a securely fastened helmet.} \par \textit {Head injuries are the leading cause of motorcycle deaths, and a properly fastened helmet is the most effective protection against them. Helmets save lives. Wear full face helmet. Don’t tempt fate. Don't test the waters.} \\
\end{minipage}

\begin{minipage}{\linewidth}
\noindent\textbf{19.}  -- \textbf {(C).} Wearing a helmet can reduce the chance of a fatal head injury in:  \textbf {Any crash, regardless of speed.} \par \textit {A helmet reduces the risk of fatal head injury in every type of crash.Helmets save lives. Wear full face helmet. Don’t tempt fate. Don't test the waters.} \\
\end{minipage}

\begin{minipage}{\linewidth}
\noindent\textbf{20.}  -- \textbf {(A).} Wearing a helmet will:  \textbf {Reduce your chances of neck and head injuries.} \par \textit {The helmet does not cause neck injure and does not restrict the view. These are all urban legends. Wear full face helmet. Don’t tempt fate. Don't test the waters.} \\
\end{minipage}

\begin{minipage}{\linewidth}
\noindent\textbf{21.}  -- \textbf {(C).} Of the following, which one gives the most eye and face protection while riding?  \textbf {A faceshield.} \par \textit {A faceshield covers the whole face and eyes against wind, debris, and impact. All other options give partial protection. Just imagine a rock flying out from under the wheels of the car ahead. Or nail. Or...} \\
\end{minipage}

\begin{minipage}{\linewidth}
\noindent\textbf{22.}  -- \textbf {(D).} While riding an unfamiliar motorcycle:  \textbf {Allow extra room for stopping.} \par \textit {Unfamiliar brakes and handling are hard to judge, so leave extra distance to stop until you get to know the bike better.} \\
\end{minipage}

\begin{minipage}{\linewidth}
\noindent\textbf{23.}  -- \textbf {(C).} The engine cut-off switch is usually located:  \textbf {Near the right hand grip.} \par \textit {The cut-off switch is a key safety control. Picture a bike sliding on its side with the rear wheel still spinning. It's usually by the right grip, so you can shut off the engine with just your thumb.} \\
\end{minipage}

\begin{minipage}{\linewidth}
\noindent\textbf{24.}  -- \textbf {(C).} With a properly adjusted mirror you can see:  \textbf {Part of the lane beside you and behind you.} \par \textit {Motorcycle mirrors show the lane behind and only part of the lane beside you, so blind spots remain and you must still turn your head to check (do a head-checks).} \\
\end{minipage}

\begin{minipage}{\linewidth}
\noindent\textbf{25.}  -- \textbf {(C).} When slowing for a normal stop:  \textbf {Shift down through the gears as you slow.} \par \textit {Downshifting as you slow keeps you in the right gear, so when you come to a full stop (like at a red light) you're already in 1st and ready to go. Using one brake alone or coasting with the clutch pulled in gives less control.} \\
\end{minipage}

\begin{minipage}{\linewidth}
\noindent\textbf{26.}  -- \textbf {(D).} When turning, you should:  \textbf {Look through the turn.} \par \textit {The bike tends to go where you look, so look through the turn to where you want to go. Fixating on the center line or guard rail pulls you toward it.} \\
\end{minipage}

\begin{minipage}{\linewidth}
\noindent\textbf{27.}  -- \textbf {(B).} Which one of the following statements is correct:  \textbf {Press left, lean left, go left.} \par \textit {Just to remember the answer: press, lean, and direction always match. The key point is to pull your chin towards the mirror. Also called as a mirror kiss. You will automatically make three things: the correct body tilt forward and to the side, press on correct grip and your head uprized. without any without any thought about left, right, push, pull, lean etc. It's magic.} \\
\end{minipage}

\begin{minipage}{\linewidth}
\noindent\textbf{28.}  -- \textbf {(A).} While in a turn, it is best to:  \textbf {Maintain a steady speed or gradually accelerate.} \par \textit {Steady or gently increasing throttle keeps the bike tilt angle stable and the tires gripping. Acceleration, slowing or braking mid-turn upsets traction, and change bike tilt angle and so the curve.} \\
\end{minipage}

\begin{minipage}{\linewidth}
\noindent\textbf{29.}  -- \textbf {(C).} If you need to slow down while turning:  \textbf {Use caution, apply both brakes.} \par \textit {Main rule is to always use both brakes. Hard braking can cause a skid, but braking is not forbidden.} \\
\end{minipage}

\begin{minipage}{\linewidth}
\noindent\textbf{30.}  -- \textbf {(B).} To make a normal stop, use:  \textbf {Both brakes, and downshift.} \par \textit {Both brakes together give the most stopping power and control, and downshifting leaves you in 1st gear ready to go. Using one brake alone reduces control.} \\
\end{minipage}

\begin{minipage}{\linewidth}
\noindent\textbf{31.}  -- \textbf {(C).} Your motorcycle has two brakes. Use both brakes:  \textbf {Every time you slow or stop.} \par \textit {The front brake provides most of the stopping power, but using both every time builds the habit you need in an emergency.} \\
\end{minipage}

\begin{minipage}{\linewidth}
\noindent\textbf{32.}  -- \textbf {(B).} In a normal turn:  \textbf {Lean your motorcycle and your body at the same angle.} \par \textit {At normal turning speeds the bike and rider lean together at the same angle, which keeps the machine balanced.} \\
\end{minipage}

\begin{minipage}{\linewidth}
\noindent\textbf{33.}  -- \textbf {(A).} In slow tight turns:  \textbf {Lean the motorcycle only, and keep your body straight.} \par \textit {At low speed the bike must lean to turn tightly, while your upright body acts as a counterweight that keeps you balanced.} \\
\end{minipage}

\begin{minipage}{\linewidth}
\noindent\textbf{34.}  -- \textbf {(D).} To reduce speed before turning:  \textbf {Close the throttle and apply both brakes.} \par \textit {Slow down before the turn. You never know how good traction will be. Build your braking habit: no throttle, both brakes.} \\
\end{minipage}

\begin{minipage}{\linewidth}
\noindent\textbf{35.}  -- \textbf {(D).} When turning, it is important to:  \textbf {Look through the turn to where you want to go.} \par \textit {The bike tends to go where you look, so look through the turn to your exit. Slowing should be done before the turn, not in or out of it. You never know how good traction you have in turn.} \\
\end{minipage}

\begin{minipage}{\linewidth}
\noindent\textbf{36.}  -- \textbf {(B).} It is important to search the road for escape routes:  \textbf {At intersections.} \par \textit {Intersections are probably the most dangerous place, so always know where you can change lanes if something does wrong.} \\
\end{minipage}

\begin{minipage}{\linewidth}
\noindent\textbf{37.}  -- \textbf {(B).} To make good judgments in traffic, you need to first:  \textbf {Search ahead.} \par \textit {Good judgment starts with information: search the road ahead first, then evaluate what you see and act on it.} \\
\end{minipage}

\begin{minipage}{\linewidth}
\noindent\textbf{38.}  -- \textbf {(B).} The car driver is signaling for a left turn. The rider should:  \textbf {Slow and move away from approaching vehicles.} \par \textit {Even with the right-of-way, a driver may misjudge your speed or not see you at all. Slowing and keeping a cushion of space around you gives you time and space to react.} \\
\end{minipage}

\begin{minipage}{\linewidth}
\noindent\textbf{39.}  -- \textbf {(D).} You should make a special point of checking traffic to the rear:  \textbf {Before slowing.} \par \textit {Check behind before slowing so you know that a following driver can see you and can stop in time. You can warn them additionaly with blinking your brake light if needed.} \\
\end{minipage}

\begin{minipage}{\linewidth}
\noindent\textbf{40.}  -- \textbf {(A).} When preparing to pass, it is usually better to:  \textbf {Signal, check your left mirror,  look left.} \par \textit {A trick question about head-checks. Signal comes first it is obvious, then mirror because it is simple, then confirmation of no-change and blind spots. Just what the doctor ordered.} \\
\end{minipage}

\begin{minipage}{\linewidth}
\noindent\textbf{41.}  -- \textbf {(C).} Riding in the center lane position:  \textbf {Keeps others from sharing your lane.} \par \textit {The center position makes you more visible and leaves no chance for another vehicle to squeeze into your lane beside you. These are a lot of urban legends about riding in the center position. Please note that although it is the best solution, it has its drawbacks.} \\
\end{minipage}

\begin{minipage}{\linewidth}
\noindent\textbf{42.}  -- \textbf {(B).} To discourage other drivers from sharing your lane, it is usually best to:  \textbf {Ride in the center portion of your lane.} \par \textit {Riding in the center takes up the lane visually, so other drivers are less tempted to pass or share it. Riding to either side invites them in. These are a lot of urban legends about riding in the center position. Please note that although it is the best solution, it has its drawbacks.} \\
\end{minipage}

\begin{minipage}{\linewidth}
\noindent\textbf{43.}  -- \textbf {(B).} Keeping a cushion of space around your motorcycle:  \textbf {Allows you time and space to react.} \par \textit {Space around you gives time to see hazards and room to brake or swerve. Every rider needs it.} \\
\end{minipage}

\begin{minipage}{\linewidth}
\noindent\textbf{44.}  -- \textbf {(A).} To spot approaching traffic and maintain a space cushion in the illustration above:  \textbf {Ride in the center of the lane.} \par \textit {Riding through winding road curve, it is dangerous to get too close to the dividing line, as oncoming traffic may cut corners and enter your lane. The right side provides a poor view of the distance ahead. The center of the lane gives room on both sides, so you keep a cushion from oncoming traffic and from the edge while you watch the road ahead.} \\
\end{minipage}

\begin{minipage}{\linewidth}
\noindent\textbf{45.}  -- \textbf {(B).} You plan to make a right turn after stopping at an intersection. Usually it is best to:  \textbf {Turn from the center portion of the lane.} \par \textit {Staying in the center portion keeps other vehicles from squeezing past you in your lane while you turn. These are a lot of urban legends about riding in the center position. Please note that although it is the best solution, it has its drawbacks.} \\
\end{minipage}

\begin{minipage}{\linewidth}
\noindent\textbf{46.}  -- \textbf {(D).} Following close to vehicles in front of you:  \textbf {Is a factor in crashes caused by motorcyclists.} \par \textit {Tailgating leaves no time or space to react or stop, and it blocks your view of the road and pavement ahead.} \\
\end{minipage}

\begin{minipage}{\linewidth}
\noindent\textbf{47.}  -- \textbf {(A).} When following a car, you should try to:  \textbf {Ride where the driver can see you in the rear-view mirror.} \par \textit {The drivers mostly check rear-view mirror only. So center portion is your best choice.} \\
\end{minipage}

\begin{minipage}{\linewidth}
\noindent\textbf{48.}  -- \textbf {(A).} At intersections, the most common cause of motorcycle/car crashes is:  \textbf {Drivers entering a rider's right of way.} \par \textit {The most common crash type is a driver failing to see the motorcycle and pulling or turning into the rider's path.} \\
\end{minipage}

\begin{minipage}{\linewidth}
\noindent\textbf{49.}  -- \textbf {(A).} Car/Motorcycle crashes are most common:  \textbf {At intersections.} \par \textit {Intersections are where paths cross and drivers most often fail to see riders.} \\
\end{minipage}

\begin{minipage}{\linewidth}
\noindent\textbf{50.}  -- \textbf {(C).} Making eye contact with a driver at an intersection:  \textbf {Does not mean that the driver sees you.} \par \textit {A driver can look right at you without registering a motorcycle, so eye contact is no guarantee of being seen or yielded to. Stay ready to react.} \\
\end{minipage}

\begin{minipage}{\linewidth}
\noindent\textbf{51.}  -- \textbf {(B).} The greatest danger for a rider when passing parked cars is:  \textbf {A car pulling out.} \par \textit {A driver pulling away from the curb may block entire stree. So you will have no escape route. This is the greatest danger, so stay alert and watch for drivers in parked cars.} \\
\end{minipage}

\begin{minipage}{\linewidth}
\noindent\textbf{52.}  -- \textbf {(D).} In this picture, the light most likely to make the motorcycle noticeable to the car driver is the:  \textbf {Turn signal.} \par \textit {The only useful light in this scenario is turning light. But be sure thery won't see it anyway.} \\
\end{minipage}

\begin{minipage}{\linewidth}
\noindent\textbf{53.}  -- \textbf {(D).} When riding in traffic at night, the best way to locate bumps in the road is to:  \textbf {Watch the taillights of the car ahead.} \par \textit {Taillights that bounce up and down show bumps or rough pavement ahead. High beams don't help much in traffic and can dazzle other drivers.} \\
\end{minipage}

\begin{minipage}{\linewidth}
\noindent\textbf{54.}  -- \textbf {(D).} The best way to help others see your motorcycle is to:  \textbf {Ride with your headlight turned on.} \par \textit {A lit headlight makes you more noticeable to other drivers, day or night. Lane position and helmet color help far less.} \\
\end{minipage}

\begin{minipage}{\linewidth}
\noindent\textbf{55.}  -- \textbf {(A).} Not turning off a turn signal is most likely to:  \textbf {Encourage other drivers to pull into your path.} \par \textit {Bike turn signal does not cancel itselft after turn. A signal left on falsely tells others you are about to turn, so drivers may pull out or turn into your path.} \\
\end{minipage}

\begin{minipage}{\linewidth}
\noindent\textbf{56.}  -- \textbf {(D).} Make a special point to use your mirrors:  \textbf {Before slowing or stopping.} \par \textit {A motorcycle is narrow, and drivers behind may not notice it slowing down. Check your mirrors first to see how close the traffic behind you is, then flash your brake light to get their attention before you slow or stop.} \\
\end{minipage}

\begin{minipage}{\linewidth}
\noindent\textbf{57.}  -- \textbf {(A).} Before changing lanes to the left:  \textbf {Check your left mirror; and look to the left.} \par \textit {Mirrors have the blind spots. Head-checks is mandatory practice — turn your head to look that way before moving over.} \\
\end{minipage}

\begin{minipage}{\linewidth}
\noindent\textbf{58.}  -- \textbf {(C).} The best gear to wear to be seen at night is:  \textbf {Reflective clothing.} \par \textit {Reflective material bounces headlight beams back to drivers. Best choice at night.} \\
\end{minipage}

\begin{minipage}{\linewidth}
\noindent\textbf{59.}  -- \textbf {(B).} To see better when following a car at night:  \textbf {Use the headlights of the vehicle ahead.} \par \textit {The headlights of the vehicle ahead more powerfull and light the road for you, so you can see farther and spot hazards.} \\
\end{minipage}

\begin{minipage}{\linewidth}
\noindent\textbf{60.}  -- \textbf {(C).} When riding at night:  \textbf {Wear reflective clothing.} \par \textit {Reflective clothing makes you much easier to see in the dark.} \\
\end{minipage}

\begin{minipage}{\linewidth}
\noindent\textbf{61.}  -- \textbf {(D).} If you lock the front tire when stopping:  \textbf {Release and reapply the front brake.} \par \textit {A trick question. Front and rear brakes follow different rules. A locked front wheel is dangerous: you lose steering and balance, and the bike can fall. Release the front brake immediately to regain traction, then reapply it with less pressure.} \\
\end{minipage}

\begin{minipage}{\linewidth}
\noindent\textbf{62.}  -- \textbf {(D).} In the picture above, the rear wheel locks while trying to stop quickly. It is usually best to:  \textbf {Keep the rear wheel locked until stopped.} \par \textit {Front and rear brakes follow different rules. A locked rear wheel should stay locked until you stop. High risk to catch highkick. This a base rule. During BRC there is an exercise how to safely unlock it without risk.} \\
\end{minipage}

\begin{minipage}{\linewidth}
\noindent\textbf{63.}  -- \textbf {(B).} When swerving, it is important to:  \textbf {Separate braking from swerving.} \par \textit {Braking and swerving both use tire traction, and doing them together can cause a skid. Swerve first, then brake once the bike is upright, or brake first and then release the brake to swerve.} \\
\end{minipage}

\begin{minipage}{\linewidth}
\noindent\textbf{64.}  -- \textbf {(A).} To ride over metal bridge gratings:  \textbf {Ride straight across.} \par \textit {Gratings make the bike feel unsteady, but the tires still grip. Stay relaxed, keep a steady speed and go straight. Weaving or angling across makes the unsteadiness worse.} \\
\end{minipage}

\begin{minipage}{\linewidth}
\noindent\textbf{65.}  -- \textbf {(B).} To ride over an object on the road:  \textbf {Rise slightly off the seat.} \par \textit {Rising slightly off the seat moves your weight onto the footpegs, which is lower on the bike and improves balance. Your bent legs also absorb the impact instead of your spine.} \\
\end{minipage}

\begin{minipage}{\linewidth}
\noindent\textbf{66.}  -- \textbf {(B).} When riding over an obstacle, it is usually best to:  \textbf {Rise slightly off the seat.} \par \textit {Standing slightly on the footpegs with your knees bent lets the bike move under you. This keeps your weight low, absorbs the shock of the obstacle, and helps you stay in control.} \\
\end{minipage}

\begin{minipage}{\linewidth}
\noindent\textbf{67.}  -- \textbf {(C).} On a paved, two-lane road, sand and gravel are most likely to collect:  \textbf {Near the sides of the road.} \par \textit {Traffic pushes loose sand and gravel out of the wheel tracks toward the edges of the road, so it builds up near the sides.} \\
\end{minipage}

\begin{minipage}{\linewidth}
\noindent\textbf{68.}  -- \textbf {(D).} On a wet road, it will usually be most slippery:  \textbf {In the center portion of the lane.} \par \textit {Normally the center of the lane is the best spot, but not in rain. Oil and fluid drips there and mixes with water, so the center gets very slick. Ride in the tire tracks instead.} \\
\end{minipage}

\begin{minipage}{\linewidth}
\noindent\textbf{69.}  -- \textbf {(B).} To handle a slippery surface, it is best to:  \textbf {Slow down before reaching it.} \par \textit {Braking, accelerating and downshifting all stress traction. Reduce speed before you reach the slippery spot, then keep the bike upright and steady while crossing.} \\
\end{minipage}

\begin{minipage}{\linewidth}
\noindent\textbf{70.}  -- \textbf {(D).} To stop on wet pavement:  \textbf {Apply the front and rear brakes together.} \par \textit {Using both brakes gives the best stopping power, even in the wet. Apply them gently and smoothly, because wet pavement locks a wheel more easily.} \\
\end{minipage}

\begin{minipage}{\linewidth}
\noindent\textbf{71.}  -- \textbf {(D).} When riding where sand and gravel have collected on paved roads, you should:  \textbf {Avoid sudden changes in speed or direction.} \par \textit {Loose material gives little grip. Sudden braking, acceleration or steering can make a tire slip, so keep your inputs smooth and gradual.} \\
\end{minipage}

\begin{minipage}{\linewidth}
\noindent\textbf{72.}  -- \textbf {(C).} If the rear tire goes flat:  \textbf {The back of the motorcycle will jerk from side to side.} \par \textit {A flat rear tire makes the back of the bike sway or jerk side to side. Steering feels heavy only with a flat front tire. You can still keep control by holding the grips firmly.} \\
\end{minipage}

\begin{minipage}{\linewidth}
\noindent\textbf{73.}  -- \textbf {(B).} If you have a flat tire while riding, hold the handle grips firmly and:  \textbf {Use the brake of the good tire or avoid braking.} \par \textit {Braking the flat tire can cause a skid and a fall. If you must brake, use the good tire's brake gently, or avoid braking and let the bike slow down.} \\
\end{minipage}

\begin{minipage}{\linewidth}
\noindent\textbf{74.}  -- \textbf {(C).} When the front tire goes flat while riding:  \textbf {Apply rear brake and steer to the edge of the road.} \par \textit {The rear tire is the good one, so use its brake gently if you must brake. Never brake the flat tire. Slow down, only then move to the edge of the road.} \\
\end{minipage}

\begin{minipage}{\linewidth}
\noindent\textbf{75.}  -- \textbf {(A).} When the front tire goes flat:  \textbf {The steering will feel heavy.} \par \textit {A flat front tire makes the steering feel heavy and sluggish. Side-to-side jerking of the rear is the sign of a flat rear tire.} \\
\end{minipage}

\begin{minipage}{\linewidth}
\noindent\textbf{76.}  -- \textbf {(A).} If your throttle sticks:  \textbf {Squeeze the clutch and use engine cut-off switch.} \par \textit {The clutch stops the bike from accelerating, but a stuck throttle would still over-rev the engine and could cause damage. The cut-off switch shuts the engine off so you can brake and pull over safely.} \\
\end{minipage}

\begin{minipage}{\linewidth}
\noindent\textbf{77.}  -- \textbf {(D).} In the picture above, it is usually best to:  \textbf {Ride in position 2 or 3.} \par \textit {Positions 2 and 3 keep you away from the hazard on the left side of the lane and leave you more room to react.} \\
\end{minipage}

\begin{minipage}{\linewidth}
\noindent\textbf{78.}  -- \textbf {(D).} In this picture, A, B, and C are riding as a group. To pass, rider B should:  \textbf {Return to the right portion of the lane behind A.} \par \textit {Riders in a group pass one at a time. Each rider returns to their original position in a formation. The idea is to restore formation in front of a passig car.} \\
\end{minipage}

\begin{minipage}{\linewidth}
\noindent\textbf{79.}  -- \textbf {(B).} To create more space in the situation pictured above, you should:  \textbf {Adjust speed and stay in the center of the lane.} \par \textit {A merging driver will likely not notice you, so give them room. Yes, you have the right-of-way, but your safety is more important.} \\
\end{minipage}

\begin{minipage}{\linewidth}
\noindent\textbf{80.}  -- \textbf {(A).} Most motorcycle/car crashes:  \textbf {Occur because a car driver pulls into the path of the motorcycle.} \par \textit {A motorcycle is narrow, and drivers may not notice it. Most crashes happen when a driver fails to see the motorcycle and turns or pulls into its path.} \\
\end{minipage}

\begin{minipage}{\linewidth}
\noindent\textbf{81.}  -- \textbf {(D).} You are preparing to move to the left lane, as pictured above. To spot cars passing beside you, it is best to:  \textbf {Turn your head to the left.} \par \textit {Mirrors don't show the blind spot, where a passing car can hide. Do the head-checks, look over your shoulder before moving over.} \\
\end{minipage}

\begin{minipage}{\linewidth}
\noindent\textbf{82.}  -- \textbf {(B).} To adjust for traffic and avoid debris in the curve pictured, ride:  \textbf {In position 2.} \par \textit {The center portion keeps you away from oncoming traffic on the left and from sand and debris near the edge on the right.} \\
\end{minipage}

\begin{minipage}{\linewidth}
\noindent\textbf{83.}  -- \textbf {(D).} To cross tracks that run parallel to your course:  \textbf {Move away, then turn and cross at an angle of at least 45 degrees.} \par \textit {Crossing tracks is dangerous. Edging across gradually can catch your tire in the groove and throw you off. Do it at 90 to 45 degrees angle. Do not enter oncoming traffic lane or shoulder.} \\
\end{minipage}

\begin{minipage}{\linewidth}
\noindent\textbf{84.}  -- \textbf {(A).} In this picture, the rider in the most dangerous position is:  \textbf {Rider B.} \par \textit {Rider B is alongside the car in the passing lane, which is the most exposed spot. B is in the car's blind spot and has the least room to react.} \\
\end{minipage}

